\documentclass[10pt,a4paper]{article}
\usepackage[utf8]{inputenc}
\usepackage[a4paper,left=3cm,right=2cm]{geometry}
\usepackage{amsmath}
\usepackage[thmmarks, amsmath, thref]{ntheorem}
\usepackage{amsfonts}
\usepackage{amssymb}
\usepackage{fancyhdr}
\pagestyle{fancy}
\usepackage{anysize}
\usepackage{graphicx}
\usepackage{float}
\usepackage{hyperref}
\usepackage{multirow}
\usepackage{enumitem}
\newtheorem{theorem}{Theorem}
\theoremsymbol{\ensuremath{\blacksquare}}
\newtheorem*{solution}{Solución:}
\usepackage{listings}
\usepackage{xcolor}

% Margenes y formalidades

\textwidth 6.25in % Margen Derecho y ancho del texto
\textheight 8.5in
\oddsidemargin 0in % Margen Izquierdo
\headheight 0in % Largo del texto
\leftmargin 1in
\parindent 0em % Salto entre la indentación (Sangría)https://www.overleaf.com/4519127669sdpptkhvsbfn
\topmargin -0.5in % Margen Superior
\parskip 0.5ex % Salto entre parrafos
\baselineskip 1pt

%\lhead{\includegraphics[scale=0.2]{logo.png}} % texto izquierda de la cabecera
%\chead{}                                      % texto centro de la cabecera
%\rhead{\includegraphics[scale=0.2]{FIG/logo-mba-2016.png}} % número de página a la derecha

\renewcommand{\headrulewidth}{0.4pt} % grosor de la línea de la cabecera
\renewcommand{\footrulewidth}{0.4pt} % grosor de la línea del pie

\begin{document}



\pagebreak
\vspace*{4mm} 

\begin{center}
\begin{huge}
\textbf{\\Ejercicios Econometría 4}
\end{Large}\\
%{\large 2023-2}\\
\end{center}

\begin{enumerate}

\item En una prueba estandarizada para postular a un cargo en especifico, se pudo encontrar, a través de una encuesta, las horas de estudio para esta prueba y la motivación por el cargo. Las notas, las horas de estudio y la motivación de 5 individuos están en la siguiente tabla:
  \[
    \begin{array}{|c|c|c|c|}
    \hline
      i & Notas (Y_i) & Horas (X_{1i}) & Motivación (X_{2i})\\\hline
      1 & 30 & 4  & 4 \\
      2 & 34 & 4  & 5 \\
      3 & 40 & 6  & 7 \\
      4 & 46 & 8  & 9 \\
      5 & 52 & 10 & 9\\
      \hline
    \end{array}
  \]

  \begin{enumerate}
    \item Estime el modelo \(Y_i=\beta_0+\beta_1X_{1i}+u_i\) e informe
      \(\hat{\beta}_0,\hat{\beta}_1, R^{2}, R^{2}_{\text{aj}}\) y
      \(\operatorname{se}(\hat{\beta}_1)\).

    \item ¿Cuándo omitir \(X_2\) sesga \(\hat{\beta}_1\)?

    \item Estime \(Y_i=\beta_0+\beta_1X_{1i}+\beta_2X_{2i}+u_i\) e informe
      \(\hat{\beta}_0,\hat{\beta}_1,\hat{\beta}_2, R^{2}, R^{2}_{\text{aj}}\) y
      sus errores estándar, sabiendo que:
      \[
(X'X)^{-1}=
\begin{pmatrix}
2.7 & 0.3 & -0.65 \\
0.3 & \;\;0.325 & \;\;-0.35 \\
-0.65 & \;\;.0.35 & \;\;0.425 \\
\end{pmatrix}
\]


    \item Compare los coeficientes, \(R^{2}\) y errores estándar de los modelos
      (a) y (c). Interprete.
  \end{enumerate}
\end{enumerate}


\section{Formulario}


\begin{itemize}

\item  Estimadores MCO para RLS.
{
$$\hat{\beta}_0 = \bar{y} - \hat{\beta}_1 \bar{x} $$
$$\hat{\beta}_1 = \frac{\sum_{i=1}^{n}(x_i - \bar{x})(y_i-\bar{y})}{\sum_{i=1}^{n}(x_i - \bar{x})^2}=\frac{Cov(x,y)}{Var(x)}$$}


\item  Estimadores MCO para RLM.
{
$${\hat{\beta}} = (X^TX)^{-1} X^TY $$}


    \item \[ R^2  = 1 - \frac{SSR}{SST} \] 
\item  $$SST = \sum_{i=1}^{n} (y_i - \bar{y})^2 \quad \quad SSR = \sum_{i=1}^{n}  \hat{u}_i^2 $$
    \item \[ \bar{R}^2 =  1 - \frac{(n-1)SSR}{(n-k-1)SST} \]

    \item  $$\hat{\sigma}^2 = \frac{\sum_{i=1}^{n} \hat{u}_i^2 }{n-k-1}= \frac{SSR}{n - k - 1}$$
    \item \[
\text{Var}(\hat{\beta}|X) = \hat{\sigma}^2(X^{\top}X)^{-1}
\]
    \end{itemize}



\end{document}